\documentclass[10pt,a4paper]{amsart}
\usepackage[margin=19mm]{geometry}
\usepackage{amsmath,amssymb,mathtools}
\usepackage[scr=rsfs,cal=euler]{mathalfa}
\usepackage{xfrac}
\usepackage[unicode,hidelinks,bookmarks=false]{hyperref}
\newcommand{\ie}{i.e.\ }
\newcommand{\cf}{cf.\ }
\newcommand{\resp}{respectively\ }
\newcommand{\Subsection}{\subsection}
\providecommand{\nobreakdash}{\nobreak\hbox{-}}
\setlength{\emergencystretch}{2em}
\multlinegap0pt
\usepackage{amssymb}
\usepackage{breqn}
\usepackage{tikz}
\usepackage{float}
\usepackage{xfrac}
\newtheorem{cor}{Corollary}
\def\C{\mathbb{C}}
\title{Asymptotics for resolutions and smoothings of Calabi-Yau conifolds}
\author{Abdou Oussama Benabida}
\date{Source: arXiv:2503.16702v3 (CC BY 4.0)}
\begin{document}
\maketitle

\noindent\emph{Source and licence.} Asymptotics for resolutions and smoothings of Calabi-Yau conifolds. Version arXiv:2503.16702v3 (CC BY 4.0), distributed by arXiv under the Creative Commons Attribution 4.0 International licence, http://creativecommons.org/licenses/by/4.0/, as recorded in the arXiv licence metadata for this exact version and shown on the abstract page https://arxiv.org/abs/2503.16702v3. Full licence notice: \texttt{licenses/arxiv-2503.16702-CC-BY-4.0.txt}.\par\smallskip

\noindent\textit{Editorial scope.} Counted unit: the article's cor cor (source0.tex lines 1599--1604). The article's complete original proof of it is delivered with it (source0.tex lines 1605--1612, one \texttt{proof} environment of 8 lines). No mathematics is written, repaired or re-derived: the statement and the proof are the article's own lines, in the article's own notation, and no retained line is substituted or re-flowed.  No further source-local context is needed: every cross-reference of every retained line resolves inside the two delivered spans.  Nothing was omitted from any retained span, so the package carries no omission record.  No retained line contains a citation, so the wrapper carries no bibliography and no bibliography entry.  No retained line contains a figure environment, an image-inclusion command or drawing code, so no image is delivered and the package declares no asset dependency.  Editorial formatting only, in addition: an AMS \texttt{amsart} preamble that restates the article's own declarations and macros used by the retained lines, namely the environments \texttt{cor} on separate counters, together with the standard packages those lines need.  The retained environments print in this document as Corollary 1 in order of appearance, whereas the article prints the same items with the numbering of its own full document.  The complete native source of the article as distributed in the arXiv e-print for this exact version is bundled unchanged as \texttt{sources/175139/source0.tex}.
\begin{cor}
    There exists a polyhomogeneous family of non-vanishing holomorphic $(n,0)$-form $\Omega_s$ on $X_s$, for $s>0$ sufficiently small, such that the polyhomogeneous family of Ricci-flat Kähler metrics $\omega_{CY,s}$ satisfies
    $$
    \omega_{CY,s} = i^{n^2} \Omega_s \wedge \overline{\Omega_s}.
    $$
\end{cor}

\begin{proof}
    Let $\tilde \Omega$ be any nowhere vanishing relative holomorphic $(n,0)$-form on the smoothing $\mathcal{X}$. Therefore, its restriction to the path $\mathcal{X}_0$ is also holomorphic and therefore, naturally extends to a polyhomogeneous section of $\left(\bigwedge^{n} \hspace{0.07cm} (^{c, \varepsilon} \hspace{0.01cm} T^{*} \mathcal{X}_b) \right)^{\C}$. Moreover, since $\omega_{CY,s}$ is a Ricci-flat Kähler metric, we get that 
    $$
    \omega_{CY,s}^n = c_s i^{n^2} \tilde \Omega_s \wedge \overline{\tilde \Omega_s}, \hspace{0.1cm} \hspace{0.1cm} c_s = \frac{\int_{X_s} \omega_s^n}{\int_{X_s} i^{n^2} \tilde \Omega_s \wedge \overline{\tilde \Omega_s}}. 
    $$
    By the Melrose pushforward theorem, we get that both $\int_{X_s} \omega_s^n$ and $\int_{X_s} i^{n^2} \tilde \Omega_s \wedge \overline{\tilde \Omega_s}$ are polyhomogeneous. Moreover, $ \int_{X_s} i^{n^2} \tilde \Omega_s \wedge \overline{\tilde \Omega_s}$ is bounded below and above by certain uniform positive numbers $M, M'>0$ near $s=0$ and similarly for $\int_{X_s} \omega_s^n$, therefore, using the Taylor expansion of $\frac{1}{1+ y}$ and $\sqrt{1 + y}$, we can conclude that $\sqrt{c}$ is polyhomogeneous. Therefore, $\Omega_s := \sqrt{c_s} \tilde \Omega_s$ satisfies the wanted properties.  
 
\end{proof}

\end{document}
